\chapter{Transformations of Random Variables}

The previous chapters studied random variables coupled probabilistically.
Now consider deterministic transformations: if $X$ has a known distribution and
$g$ is a function, what is the distribution of the transformed variable $Y = g(X)$?

This question arises whenever we measure one quantity but need another:
computing dissipated power $X^2$ from voltage $X$, or converting uniform random
numbers into Gaussian variables. We develop the cumulative distribution function
(CDF) method, the Jacobian change-of-variables formula for smooth maps, and the
distributions of order statistics.

\chapterline

\section{The Core Question, and Why It Is Subtle}
\label{ch16:core}

Let $X$ be a random variable with a known distribution, and let $g$ be a function. Then
$Y = g(X)$ is a new random variable. Its randomness is entirely inherited from $X$ -- once
the universe picks the value of $X$, the value of $Y$ is determined. So in some sense
nothing new is happening. And yet the \emph{distribution} of $Y$ can look wildly different
from that of $X$.

\begin{intuitionbox}
Picture the number line for $X$, and imagine its probability spread along it like a thin
layer of paint -- thick where the density is high, thin where it is low. The function $g$
\emph{relocates} every point: the paint at position $x$ gets moved to position $g(x)$.
If $g$ stretches a region, the paint there gets spread thinner; if $g$ compresses a
region, the paint piles up. The distribution of $Y$ is just the picture of where all the
paint ended up. The whole subject is bookkeeping for \emph{conservation of paint}:
probability is never created or destroyed, only moved.
\end{intuitionbox}

Here is the subtlety that trips everyone the first time. It is tempting to guess that the
density of $Y$ at $y$ is simply the density of $X$ at $x = g^{-1}(y)$. That is wrong, and
understanding why is what this chapter is really about.

\begin{trapbox}
\textbf{The most common mistake.} If $Y = g(X)$, then \emph{it is not true} that
$f_Y(y) = f_X\!\big(g^{-1}(y)\big)$. Densities are not heights you can just relabel; they
are \emph{amounts of probability per unit length}. When $g$ stretches the axis, the same
chunk of probability is smeared over a longer interval, so the density must \emph{drop}.
You must account for how much $g$ stretches or compresses -- that correction factor is the
Jacobian. Forgetting it is the number-one error in transformation problems.

To see how badly naive relabeling fails, consider $Y = 2X$ with $X \sim \Unif(0,1)$. The
inverse map is $x = y/2$. If you simply set $f_Y(y) = f_X(y/2) = 1$ on $(0,2)$, the total
probability would be $\int_0^2 1\,\dd y = 2$ -- twice the total probability of the entire
sample space. Because the interval was stretched by a factor of $2$, the density must be
halved to $f_Y(y) = 1/2$.
\end{trapbox}

A clean way to never make that mistake is to always reason in terms of the CDF, which
deals in \emph{accumulated probability} (a dimensionless quantity that is immune to
stretching) rather than density. So we start there.

\section{The CDF Method}
\label{ch16:cdf-method}

The cumulative distribution function $F_Y(y) = \P(Y \le y)$ is the most robust handle on a
distribution, because it is defined for \emph{every} random variable -- discrete,
continuous, mixed, weird -- and because it is just a probability, so we can manipulate it
with the laws of probability we already trust.

\begin{focusbox}
\textbf{The CDF recipe.} To find the distribution of $Y = g(X)$:
\begin{enumerate}
  \item Write $F_Y(y) = \P(Y \le y) = \P\big(g(X) \le y\big)$.
  \item \emph{Solve the inequality} $g(X) \le y$ for $X$, turning it into a statement
        about $X$ landing in some set.
  \item Evaluate that probability using the known distribution of $X$ (i.e.\ using
        $F_X$).
  \item If $Y$ is continuous, differentiate: $f_Y(y) = F_Y'(y)$.
\end{enumerate}
The only step that takes thought is step 2 -- and that is just algebra plus care about the
direction of inequalities.
\end{focusbox}

\subsection{First worked example: the square of a uniform}

\begin{example}[$Y=X^2$ for a uniform]
\label{ch16:unif-square}
Let $X \sim \Unif(0,1)$, so $f_X(x) = 1$ for $0 < x < 1$, and let $Y = X^2$. Since $X$
lives in $(0,1)$, so does $Y$. For $0 < y < 1$,
\[
  F_Y(y) = \P(X^2 \le y) = \P(X \le \sqrt{y}) = F_X(\sqrt{y}) = \sqrt{y},
\]
where we used that $X \ge 0$ so squaring is increasing on the relevant range.
Differentiating,
\[
  f_Y(y) = \frac{\dd}{\dd y}\sqrt{y} = \frac{1}{2\sqrt{y}}, \qquad 0 < y < 1.
\]
\end{example}

Stop and look at what just happened. We started with a perfectly flat density (the
uniform) and ended with one that \emph{blows up} near $0$ (Figure~\ref{ch16:fig-square}).
Squaring crushes the interval $(0,1)$ toward $0$ -- the values near $1$ stay near $1$, but
the values near $\tfrac12$ get pulled down to $\tfrac14$, and values near $0.1$ get pulled
all the way to $0.01$. So probability piles up near the origin, and the density spikes
there. The CDF method delivered this automatically; we never had to ``think'' about
paint piling up, because $F$ does that accounting for free.

\begin{figure}[H]
\centering
\begin{tikzpicture}
\begin{axis}[
  width=11.5cm, height=6.2cm,
  xlabel={value}, ylabel={density},
  domain=0.02:1, samples=120,
  ymin=0, ymax=4, xmin=0, xmax=1,
  legend pos=north east, legend cell align=left,
  axis lines=left, thick,
]
\addplot[densityred, very thick] {1/(2*sqrt(x))};
\addlegendentry{$f_Y(y)=\tfrac{1}{2\sqrt{y}}$, \ $Y=X^2$}
\addplot[distblue, very thick, dashed] coordinates {(0.02,1)(1,1)};
\addlegendentry{$f_X(x)=1$, \ $X\sim\Unif(0,1)$}
\end{axis}
\end{tikzpicture}
\caption{Squaring a uniform. The flat density of $X$ (dashed) becomes the spiking density
of $Y=X^2$ (solid). Probability is conserved -- both curves enclose area $1$ -- but
squaring compresses the small values together, piling density up near $0$.}
\label{ch16:fig-square}
\end{figure}

\begin{recoverybox}
\textbf{Pause and think.} Suppose $X \sim \Unif(0, 1)$ and $Y = 2X$. Before applying any
formula, where does $Y$ live, and what must its density be if total probability is conserved?
\emph{Check:} $Y$ lives on $(0, 2)$, an interval of width $2$. Conserving the total probability
of $1$ forces the density to be constant and equal to $1/2$. Doubling the interval width halves
the density height.
\end{recoverybox}

\subsection{When $g$ is not monotone: fold the probability}

If $g$ is not one-to-one, the inequality $g(X)\le y$ can correspond to \emph{several}
intervals of $X$, and the CDF method handles this with no extra machinery -- you just add
up all the pieces.

\begin{example}[$Y = X^2$ for a standard normal -- the chi-squared]
\label{ch16:normal-square}
Let $Z \sim \Normal(0,1)$ and $Y = Z^2$. Now $Z$ ranges over all of $\R$, and the event
$\{Z^2 \le y\}$ for $y > 0$ is the event $\{-\sqrt{y} \le Z \le \sqrt{y}\}$ -- a single
interval, but a \emph{two-sided} one. Thus
\[
  F_Y(y) = \P(-\sqrt{y} \le Z \le \sqrt{y}) = \Phi(\sqrt{y}) - \Phi(-\sqrt{y})
         = 2\Phi(\sqrt{y}) - 1,
\]
using the symmetry $\Phi(-a) = 1 - \Phi(a)$. Differentiate, recalling
$\Phi'(z) = \varphi(z) = \tfrac{1}{\sqrt{2\pi}}e^{-z^2/2}$ and the chain rule
$\tfrac{\dd}{\dd y}\sqrt{y} = \tfrac{1}{2\sqrt{y}}$:
\[
  f_Y(y) = 2\,\varphi(\sqrt{y})\cdot \frac{1}{2\sqrt{y}}
         = \frac{1}{\sqrt{2\pi}}\, y^{-1/2} e^{-y/2}, \qquad y > 0.
\]
This is the \textbf{chi-squared distribution with one degree of freedom}, denoted $\chi^2_1$.
It measures the squared deviation or kinetic energy of a single standard Gaussian kick. (The
term ``degree of freedom'' simply records how many independent squared standard normals are
being added together -- here, exactly one.) It is also a member of the Gamma family,
$\Gam(\tfrac12,\tfrac12)$, which governs waiting times and squared Gaussian magnitudes.
The factor of $2$ came from \emph{both} tails of the normal mapping onto the same positive
$y$ -- the probability got \emph{folded}.
\end{example}

\begin{whybox}
Why the factor of $2$? Because $g(z)=z^2$ maps two distinct inputs, $+\sqrt{y}$ and
$-\sqrt{y}$, to the same output $y$. Each contributes its own sliver of probability to the
density at $y$. The general lesson: when $g$ is many-to-one, the density of $Y$ at $y$ is a
\emph{sum} over all the preimages $x$ with $g(x)=y$. The CDF method discovers this sum on
its own, simply because solving $g(X)\le y$ naturally collects every relevant region of
$X$.
\end{whybox}

\subsection{The general monotone formula, derived}

Let us extract the pattern. Suppose $g$ is \emph{strictly increasing} and differentiable
on the range of $X$, with inverse $g^{-1}$. Then $g(X)\le y \iff X \le g^{-1}(y)$, so
\[
  F_Y(y) = \P\big(X \le g^{-1}(y)\big) = F_X\big(g^{-1}(y)\big).
\]
Differentiating with the chain rule,
\[
  f_Y(y) = f_X\big(g^{-1}(y)\big)\cdot \frac{\dd}{\dd y} g^{-1}(y).
\]
If instead $g$ is \emph{strictly decreasing}, then $g(X)\le y \iff X \ge g^{-1}(y)$, so
$F_Y(y) = 1 - F_X(g^{-1}(y))$. Differentiating yields
\[
  f_Y(y) = \frac{\dd}{\dd y}\big(1 - F_X(g^{-1}(y))\big)
         = -f_X(g^{-1}(y))\cdot \frac{\dd}{\dd y}g^{-1}(y)
         = f_X(g^{-1}(y))\cdot \left(-\frac{\dd}{\dd y}g^{-1}(y)\right).
\]
Because $g$ is decreasing, its inverse $g^{-1}$ is also decreasing, meaning
$\frac{\dd}{\dd y}g^{-1}(y) < 0$. The factor $-\frac{\dd}{\dd y}g^{-1}(y)$ is therefore
strictly positive and equals $\abs{\frac{\dd}{\dd y}g^{-1}(y)}$. Both cases collapse into
one clean statement, which is important enough to box.

\begin{formalbox}
\textbf{Change-of-variables formula (monotone case).} If $g$ is strictly monotone and
differentiable with differentiable inverse, and $X$ has density $f_X$, then $Y=g(X)$ has
density
\[
  \boxed{\;f_Y(y) = f_X\big(g^{-1}(y)\big)\,\abs[\Big]{\frac{\dd}{\dd y}\,g^{-1}(y)}\;}
\]
on the range of $g$, and $0$ elsewhere. Writing $x = g^{-1}(y)$ and using
$\tfrac{\dd x}{\dd y} = 1/g'(x)$, this is the same as
$f_Y(y) = f_X(x)\,/\,\abs{g'(x)}$.
\end{formalbox}

That absolute-value derivative term is the \textbf{Jacobian} (in one dimension). The
Jacobian is the local stretch factor: it relates the width of an infinitesimal output
interval $\dd y$ to the width of the corresponding input interval $\dd x$, via
$\dd y = \abs{g'(x)}\,\dd x$. Where $g$ stretches space ($\abs{g'} > 1$), a fixed slice of
probability is smeared over a wider interval and the density must drop; where $g$
compresses space ($\abs{g'} < 1$), the paint piles up and the density spikes. This is the
precise, quantitative form of ``conservation of paint.''

\section{The Change-of-Variables (Jacobian) Method}
\label{ch16:jacobian}

\begin{connectionbox}
\textbf{This is just $u$-substitution.} Recall from calculus that if you substitute
$y = g(x)$ in an integral, then $\dd y = g'(x)\,\dd x$, and the limits and the integrand
both transform. The transformation formula for densities is that same substitution rule
applied to the statement $\int_A f_X(x)\,\dd x = \P(X \in A)$. We are literally
re-expressing the probability integral in the new variable $y$. The absolute value on the
Jacobian appears because probability is the \emph{unsigned} area under a density --
swapping the orientation of the limits (which a decreasing $g$ does) must not flip the sign
of a probability. See the companion calculus volume's chapter on substitution and the
change-of-variables theorem for the multivariable Jacobian determinant we use below.
\end{connectionbox}

\begin{whybox}
\textbf{The infinitesimal derivation.} Forget formulas for a moment and think in
infinitesimals. A tiny interval $[x, x+\dd x]$ carries probability $f_X(x)\,\dd x$. Under
a smooth $g$ it maps to the tiny interval $[y, y+\dd y]$ with $\dd y = \abs{g'(x)}\,\dd x$.
\emph{All} of that probability must end up in the image interval -- paint is conserved --
so
\[
  f_Y(y)\,\dd y = f_X(x)\,\dd x
  \quad\Longrightarrow\quad
  f_Y(y) = f_X(x)\,\frac{\dd x}{\dd y} = \frac{f_X(x)}{\abs{g'(x)}}.
\]
That is the whole theorem in two lines. Everything else is a careful version of this:
``equal probability in corresponding little boxes.''
\end{whybox}

\begin{figure}[H]
\centering
\begin{tikzpicture}[scale=1.05]
  \draw[->, thick, gray!70] (-0.3, 0) -- (5.2, 0) node[right, font=\small] {$x$};
  \draw[->, thick, gray!70] (0, -0.3) -- (0, 4.6) node[above, font=\small] {$y$};

  \def\xone{1.3}
  \def\xtwo{1.8}
  \pgfmathsetmacro{\yone}{0.4 + 0.45*\xone*\xone}
  \pgfmathsetmacro{\ytwo}{0.4 + 0.45*\xtwo*\xtwo}

  \fill[sourceblue!25, draw=sourceblue!90!black, thick] (\xone, 0) rectangle (\xtwo, -0.45);
  \node[font=\footnotesize, sourceblue!90!black] at ({0.5*(\xone+\xtwo)}, -0.22) {$\dd x$};
  \node[font=\scriptsize, align=center, sourceblue!90!black, below=2pt] at ({0.5*(\xone+\xtwo)}, -0.45) {$\dd P = f_X(x)\,\dd x$};

  \fill[densityred!25, draw=densityred!90!black, thick] (0, \yone) rectangle (-0.45, \ytwo);
  \node[font=\footnotesize, densityred!90!black] at (-0.22, {0.5*(\yone+\ytwo)}) {$\dd y$};
  \node[font=\scriptsize, align=right, densityred!90!black, left=3pt] at (-0.45, {0.5*(\yone+\ytwo)}) {$\dd P = f_Y(y)\,\dd y$};

  \draw[dashed, sourceblue!70, thick] (\xone, 0) -- (\xone, \yone) -- (0, \yone);
  \draw[dashed, sourceblue!70, thick] (\xtwo, 0) -- (\xtwo, \ytwo) -- (0, \ytwo);

  \draw[domain=0.3:2.8, smooth, variable=\t, very thick, targetgreen!80!black]
    plot ({\t}, {0.4 + 0.45*\t*\t}) node[right, font=\small] {$y = g(x)$};

  \pgfmathsetmacro{\xm}{0.5*(\xone+\xtwo)}
  \pgfmathsetmacro{\ym}{0.4 + 0.45*\xm*\xm}
  \pgfmathsetmacro{\slope}{0.9*\xm}
  \draw[very thick, quotientorange!90!black] ({\xm-0.6}, {\ym-0.6*\slope}) -- ({\xm+0.6}, {\ym+0.6*\slope});
  \fill[quotientorange!90!black] (\xm, \ym) circle (2pt);
  \node[font=\footnotesize\bfseries, quotientorange!90!black, anchor=south east] at ({\xm-0.1}, {\ym+0.4})
    {$\text{Slope } g'(x) = \dfrac{\dd y}{\dd x}$};

  \node[draw=gray!40, fill=gray!3, rounded corners=3pt, align=center, font=\small, inner sep=5pt]
    at (3.6, 1.8) {Conservation of mass:\\[3pt]
    $f_Y(y)\,\dd y = f_X(x)\,\dd x$\\[4pt]
    $\implies f_Y(y) = \dfrac{f_X(x)}{g'(x)}$};
\end{tikzpicture}
\caption{The geometry of the one-dimensional Jacobian. An input interval of width $\dd x$ carries probability mass $\dd P = f_X(x)\,\dd x$. The transformation $y = g(x)$ maps it to an output interval of width $\dd y = g'(x)\,\dd x$. Because probability is conserved ($\dd P = f_X(x)\,\dd x = f_Y(y)\,\dd y$), the output density must equal the input density divided by the local slope: $f_Y(y) = f_X(x)/g'(x)$. Where the curve is steep, the interval is stretched and the density is diluted.}
\label{ch16:fig-jacobian-1d}
\end{figure}

\subsection{Worked example: the lognormal}

\begin{example}[Exponentiating a normal]
\label{ch16:lognormal}
Let $X \sim \Normal(\mu,\sigma^2)$ and $Y = e^{X}$. Here $g(x) = e^x$ is strictly
increasing with inverse $g^{-1}(y) = \ln y$ for $y>0$, and $\tfrac{\dd}{\dd y}\ln y = 1/y$.
Plugging into the boxed formula,
\[
  f_Y(y) = f_X(\ln y)\cdot \frac{1}{y}
         = \frac{1}{y\,\sigma\sqrt{2\pi}}\,
           \exp\!\left(-\frac{(\ln y - \mu)^2}{2\sigma^2}\right), \qquad y>0.
\]
This is the \textbf{lognormal distribution}: a random variable whose logarithm is normal.
It describes quantities formed by compounding many independent percentage changes -- such
as stock prices, incomes, and sizes of cities -- yielding a distribution strictly confined
to positive numbers with a long, heavy right tail.
\end{example}

\begin{figure}[H]
\centering
\begin{tikzpicture}
\begin{axis}[
  width=11.5cm, height=6cm,
  xlabel={$y$}, ylabel={density},
  domain=0.02:6, samples=160,
  ymin=0, ymax=1.0, xmin=0, xmax=6,
  legend pos=north east, legend cell align=left,
  axis lines=left, thick,
]
\addplot[densityred, very thick] {1/(x*0.5*sqrt(2*pi))*exp(-(ln(x))^2/(2*0.25))};
\addlegendentry{$\sigma=0.5$}
\addplot[distblue, very thick, dashed] {1/(x*1*sqrt(2*pi))*exp(-(ln(x))^2/(2*1))};
\addlegendentry{$\sigma=1$}
\end{axis}
\end{tikzpicture}
\caption{Lognormal densities ($\mu=0$). The variable lives on $(0,\infty)$, is skewed to
the right, and has a long upper tail -- the signature of a quantity built from
multiplicative shocks. The Gaussian exponential factor pins the density to zero at the origin, overcoming the diverging $1/y$ Jacobian factor.}
\label{ch16:fig-lognormal}
\end{figure}

\begin{applicationbox}
\textbf{Why finance loves the lognormal.} Stock prices cannot go negative, and over a day
they tend to move by \emph{percentages}, not by fixed dollar amounts -- a \$1000 stock and
a \$10 stock both ``feel'' a $2\%$ move. If the log-return $\ln(S_{t+1}/S_t)$ is roughly
normal (a sum of many small log-shocks, governed by the central limit theorem),
then the price ratio is lognormal (Figure~\ref{ch16:fig-lognormal}), and the price stays positive automatically. The
celebrated Black--Scholes option-pricing model is built on exactly this transformation:
model the log-price as Brownian motion (the continuous-time mathematical limit of a particle
buffeted incessantly by tiny random shocks, producing normal increments), and the price
itself is lognormal. The $1/y$ Jacobian factor is the mathematical fingerprint of ``returns
are multiplicative.''
\end{applicationbox}

\subsection{The affine case, and why normals stay normal}

\begin{example}[Linear transformation of a normal]
\label{ch16:affine-normal}
Let $X \sim \Normal(\mu,\sigma^2)$ and $Y = aX + b$ with $a \ne 0$. Then
$g^{-1}(y) = (y-b)/a$ and $\abs{\tfrac{\dd}{\dd y}g^{-1}(y)} = 1/\abs{a}$, so
\[
  f_Y(y) = \frac{1}{\abs{a}}\,
           \frac{1}{\sigma\sqrt{2\pi}}\exp\!\left(-\frac{((y-b)/a-\mu)^2}{2\sigma^2}\right)
         = \frac{1}{\abs{a}\sigma\sqrt{2\pi}}
           \exp\!\left(-\frac{(y-(a\mu+b))^2}{2(a\sigma)^2}\right).
\]
This is exactly the $\Normal(a\mu+b,\, a^2\sigma^2)$ density. So a linear function of a
normal is normal, with the mean and variance transforming as $\E[Y]=a\mu+b$ and
$\Var(Y)=a^2\sigma^2$ -- precisely the rules you learned for $\E$ and $\Var$, now confirmed
at the level of the \emph{whole distribution}. Setting $a=1/\sigma$, $b=-\mu/\sigma$
recovers the standardization $Z=(X-\mu)/\sigma \sim \Normal(0,1)$.
\end{example}

\subsection{Several variables at once: the Jacobian determinant}

When $g$ maps a vector $(X_1,\dots,X_n)$ to a vector $(Y_1,\dots,Y_n)$, the one-dimensional
stretch factor $\abs{g'}$ generalizes to the absolute value of the \textbf{Jacobian determinant}
of the inverse map $\mathbf{x} = h(\mathbf{y})$. The Jacobian determinant measures how much an
infinitesimal $n$-dimensional block of volume in target space scales when pulled back into source
space: $\dd x_1\cdots\dd x_n = \abs{\det J_h(\mathbf y)}\,\dd y_1\cdots\dd y_n$. The logic is
unchanged -- equal probability in corresponding little boxes -- but now the ``little box'' is
an $n$-dimensional parallelotope, scaled by the linear approximation's determinant.

\begin{formalbox}
\textbf{Multivariate change of variables.} Let $\mathbf{X}=(X_1,\dots,X_n)$ have joint
density $f_{\mathbf X}$, and let $\mathbf Y = g(\mathbf X)$ for a smooth invertible $g$
with inverse $\mathbf x = h(\mathbf y)$. Then
\[
  f_{\mathbf Y}(\mathbf y)
  = f_{\mathbf X}\big(h(\mathbf y)\big)\,\abs[\big]{\det J_h(\mathbf y)},
  \qquad
  J_h(\mathbf y) = \left[\frac{\partial x_i}{\partial y_j}\right]_{i,j}.
\]
\end{formalbox}

\begin{figure}[H]
\centering
\begin{tikzpicture}[scale=0.95]
  % Panel 1: Target coordinate plane (y_1, y_2)
  \begin{scope}[shift={(0,0)}]
    \draw[draw=gray!30, fill=gray!3, rounded corners=4pt] (-0.6,-1.8) rectangle (4.8,3.6);
    \node[anchor=north, font=\small\bfseries, sourceblue!80!black] at (2.1,3.4) {1. Coordinate Box in $y$-Plane};

    % Axes
    \draw[->, thick, gray!70] (-0.2,0) -- (4.4,0) node[right, font=\scriptsize] {$y_1$};
    \draw[->, thick, gray!70] (0,-0.2) -- (0,2.3) node[above, font=\scriptsize] {$y_2$};

    % Rectangular tile at (1.3, 0.7) with dimensions (2.0, 1.3)
    \coordinate (Y0) at (1.3, 0.7);
    \fill[sourceblue!25, draw=sourceblue!90!black, thick] (Y0) rectangle ($(Y0)+(2.0, 1.3)$);

    % Dimension labels
    \draw[<->, >=stealth, gray!80] ($(Y0)+(0, -0.25)$) -- ($(Y0)+(2.0, -0.25)$)
      node[midway, below, font=\scriptsize] {$\dd y_1$};
    \draw[<->, >=stealth, gray!80] ($(Y0)+(-0.25, 0)$) -- ($(Y0)+(-0.25, 1.3)$)
      node[midway, left, font=\scriptsize] {$\dd y_2$};

    \node[font=\scriptsize\bfseries, sourceblue!90!black] at ($(Y0)+(1.0, 0.65)$) {$\text{Area} = \dd y_1 \dd y_2$};

    \node[align=center, font=\scriptsize, draw=sourceblue!60, fill=white, rounded corners=2pt, inner sep=2pt]
      at (2.1, -1.1) {Probability in Tile:\\[1pt]$\dd P = f_{\mathbf Y}(y_1, y_2)\,\dd y_1\,\dd y_2$};
  \end{scope}

  % Transmission arrow
  \draw[->, >=stealth, very thick, structurepurple] (5.0, 1.4) -- (6.7, 1.4);
  \node[above=3pt, font=\scriptsize\bfseries, align=center, structurepurple] at (5.85, 1.4) {Inverse Map\\[1pt]\scriptsize $\mathbf{x} = h(\mathbf{y})$};

  % Panel 2: Deformed parallelogram in source plane (x_1, x_2)
  \begin{scope}[shift={(7.55,0)}]
    \draw[draw=gray!30, fill=gray!3, rounded corners=4pt] (-0.6,-1.8) rectangle (5.5,3.6);
    \node[anchor=north, font=\small\bfseries, densityred!80!black] at (2.45,3.4) {2. Deformed Parallelogram in $x$-Plane};

    % Axes
    \draw[->, thick, gray!70] (-0.2,0) -- (5.1,0) node[right, font=\scriptsize] {$x_1$};
    \draw[->, thick, gray!70] (0,-0.2) -- (0,2.5) node[above, font=\scriptsize] {$x_2$};

    % Parallelogram base point x0 = h(y0)
    \coordinate (X0) at (0.75, 0.45);
    \coordinate (V1) at (3.25, 0.35);  % dx/dy1 * dy1
    \coordinate (V2) at (0.65, 1.48);  % dx/dy2 * dy2

    \fill[densityred!25, draw=densityred!90!black, thick]
      (X0) -- ($(X0)+(V1)$) -- ($(X0)+(V1)+(V2)$) -- ($(X0)+(V2)$) -- cycle;

    % Tangent vector arrows
    \draw[->, >=stealth, very thick, targetgreen!80!black] (X0) -- ($(X0)+(V1)$)
      node[below right=-2pt, font=\tiny\bfseries] {$\frac{\partial \mathbf{x}}{\partial y_1}\dd y_1$};
    \draw[->, >=stealth, very thick, quotientorange!90!black] (X0) -- ($(X0)+(V2)$)
      node[above left=-2pt, font=\tiny\bfseries] {$\frac{\partial \mathbf{x}}{\partial y_2}\dd y_2$};

    \node[font=\tiny\bfseries, densityred!90!black] at ($(X0)+(1.95, 0.91)$) {$\text{Area} = |\det J_h|\,\dd y_1 \dd y_2$};

    \node[align=center, font=\scriptsize, draw=densityred!60, fill=white, rounded corners=2pt, inner sep=2pt]
      at (2.4, -1.1) {Conservation of Probability:\\[1pt]$f_{\mathbf Y}(\mathbf y)\,\dd y_1\dd y_2 = f_{\mathbf X}(h(\mathbf y))\,|\det J_h|\,\dd y_1\dd y_2$};
  \end{scope}
\end{tikzpicture}
\caption{The geometry of the multivariable Jacobian determinant. Under the inverse transformation $\mathbf{x} = h(\mathbf{y})$, the infinitesimal rectangular tile $\dd y_1 \times \dd y_2$ pulls back to a deformed parallelogram spanned by the partial-derivative vectors $\frac{\partial \mathbf{x}}{\partial y_1}\dd y_1$ and $\frac{\partial \mathbf{x}}{\partial y_2}\dd y_2$. The area of this parallelogram is scaled by the absolute Jacobian determinant $|\det J_h|$. Conserving probability mass across the two tiles forces $f_{\mathbf Y}(\mathbf y) = f_{\mathbf X}(h(\mathbf y))\,|\det J_h|$.}
\label{ch16:fig-jacobian-tiles}
\end{figure}

\begin{connectionbox}
The determinant here is the \emph{signed volume scaling factor} of the linear map that best
approximates $h$ near a point -- the same determinant from linear algebra. Its absolute
value answers ``how much does this map blow up or shrink volumes?'' See the companion
linear algebra volume's treatment of determinants as volume and the multivariable calculus
chapter on the change-of-variables theorem; probability simply applies that theorem to a
density.
\end{connectionbox}

We will put this two-dimensional version to spectacular use in the Box--Muller derivation
of Section~\ref{ch16:box-muller}.

\section{Order Statistics}
\label{ch16:order-stats}

So far $g$ has acted on a single variable. A different and enormously useful family of
transformations acts on an entire i.i.d.\ sample by \emph{sorting} it. Suppose
$X_1,\dots,X_n$ are i.i.d.\ with CDF $F$ and density $f$. Sort them from smallest to
largest:
\[
  X_{(1)} \le X_{(2)} \le \dots \le X_{(n)}.
\]
The $k$-th value $X_{(k)}$ is the \textbf{$k$-th order statistic}. An order statistic is
simply the value occupying rank $k$ after sorting an observed sample from smallest to largest --
the sample minimum, the sample median, or the sample maximum. In particular
$X_{(1)} = \min_i X_i$ and $X_{(n)} = \max_i X_i$. These are the random variables behind
record-breaking, extreme weather, system reliability (a chain fails when its
\emph{weakest} link fails -- the minimum), the sample median, and the range of a dataset.

\subsection{Maximum and minimum: independence does the work}

\begin{intuitionbox}
The maximum is $\le y$ exactly when \emph{every single} observation is $\le y$. Because the
observations are independent, that joint event factors into a product. That one
observation -- ``all of them'' becomes ``each of them, multiplied'' -- is the entire idea
behind the extreme order statistics.
\end{intuitionbox}

For the maximum $M = X_{(n)}$:
\[
  F_M(y) = \P(\max_i X_i \le y) = \P(X_1\le y,\dots,X_n\le y)
         = \prod_{i=1}^n \P(X_i\le y) = [F(y)]^n.
\]
Differentiating, $f_M(y) = n\,[F(y)]^{n-1} f(y)$.

For the minimum $L = X_{(1)}$, it is cleaner to use the \textbf{survival function}
$S(y) = \P(X > y) = 1 - F(y)$ -- the probability that a random lifespan outlasts threshold
$y$ -- since the minimum exceeds $y$ if and only if \emph{all} observations exceed $y$:
\[
  \P(L > y) = \P(X_1 > y,\dots,X_n > y) = [1-F(y)]^n,
  \qquad F_L(y) = 1 - [1-F(y)]^n,
\]
and $f_L(y) = n\,[1-F(y)]^{n-1} f(y)$.

\begin{example}[Minimum of exponentials -- the key to Poisson processes]
\label{ch16:min-expo}
Let $X_1,\dots,X_n$ be i.i.d.\ $\Expo(\lambda)$, so $1-F(y)=e^{-\lambda y}$. Then
$\P(L>y) = (e^{-\lambda y})^n = e^{-n\lambda y}$, which is the survival function of an
$\Expo(n\lambda)$. So the minimum of $n$ independent exponentials is again exponential,
with the rates \emph{added}. This memorable fact -- ``the first of many independent alarm
clocks rings at the combined rate'' -- is the algebraic heart of the Poisson process you
will meet later.
\end{example}

\begin{recoverybox}
\textbf{Pause and think.} You have $100$ independent lightbulbs, each with lifetime
$\Expo(\lambda)$ (mean lifetime $1/\lambda$). What is the expected waiting time until the
\emph{first} bulb burns out?
\emph{Check:} By Example~\ref{ch16:min-expo}, the minimum of the $100$ lifespans is distributed
as $\Expo(100\lambda)$. Its expected value is $\frac{1}{100\lambda}$ -- exactly one-hundredth
of a single bulb's lifespan. More competing components make the first failure arrive
proportionally faster.
\end{recoverybox}

\subsection{The general $k$-th order statistic}

\begin{whybox}
\textbf{A density by counting arrangements.} What is the density of $X_{(k)}$ at $y$? For
$X_{(k)}$ to lie in the sliver $[y,y+\dd y]$, we need: exactly one observation in that
sliver, exactly $k-1$ observations below $y$, and exactly $n-k$ above $y+\dd y$. Choosing
\emph{which} observations play these roles is a multinomial count,
$\tfrac{n!}{(k-1)!\,1!\,(n-k)!}$. Each ``below'' observation
contributes probability $F(y)$, each ``above'' contributes $1-F(y)$, and the one in the
sliver contributes $f(y)\,\dd y$. Multiply:
\[
  f_{X_{(k)}}(y) = \frac{n!}{(k-1)!\,(n-k)!}\,[F(y)]^{k-1}\,[1-F(y)]^{n-k}\,f(y).
\]
\end{whybox}

\begin{formalbox}
\textbf{Density of the $k$-th order statistic.} For an i.i.d.\ sample of size $n$ from a
continuous distribution with CDF $F$ and density $f$,
\[
  f_{X_{(k)}}(y) = \frac{n!}{(k-1)!\,(n-k)!}\,[F(y)]^{k-1}\,[1-F(y)]^{n-k}\,f(y).
\]
Setting $k=1$ and $k=n$ recovers the minimum and maximum formulas above.
\end{formalbox}

\begin{example}[Order statistics of the uniform are Beta]
\label{ch16:unif-order-beta}
Let $X_1,\dots,X_n$ be i.i.d.\ $\Unif(0,1)$, so $F(y)=y$ and $f(y)=1$ on $(0,1)$. Then
\[
  f_{X_{(k)}}(y) = \frac{n!}{(k-1)!\,(n-k)!}\,y^{k-1}(1-y)^{n-k}, \qquad 0<y<1.
\]
This is exactly the $\Bet(k,\,n-k+1)$ density! The Beta distribution is the continuous family
supported on the unit interval $(0,1)$ with density proportional to $y^{\alpha-1}(1-y)^{\beta-1}$,
making it the canonical model for random proportions and sorted positions of uniform draws.
Here $\E[X_{(k)}] = k/(n+1)$: the $n$ order statistics divide $(0,1)$ into $n+1$ gaps of equal
expected width $1/(n+1)$. The sample maximum ($k=n$) is $\Bet(n,1)$ with density $n y^{n-1}$,
matching $n[F(y)]^{n-1}f(y)$ from before.
\end{example}

\begin{figure}[H]
\centering
\begin{tikzpicture}
\begin{axis}[
  width=12cm, height=6.2cm,
  xlabel={$y$}, ylabel={density of $X_{(k)}$},
  domain=0:1, samples=120,
  ymin=0, ymax=5, xmin=0, xmax=1,
  legend style={at={(0.5,0.96)}, anchor=north}, legend cell align=left,
  axis lines=left, thick,
]
\addplot[distblue, very thick] {5*(1-x)^4};
\addlegendentry{$k=1$ (min)}
\addplot[targetgreen, very thick] {30*x^2*(1-x)^2};
\addlegendentry{$k=3$ (median)}
\addplot[densityred, very thick] {5*x^4};
\addlegendentry{$k=5$ (max)}
\end{axis}
\end{tikzpicture}
\caption{Order statistics of $n=5$ i.i.d.\ uniforms, each a Beta density. The minimum
(left) clusters near $0$, the maximum (right) near $1$, and the median (middle) peaks at
$\tfrac12$ -- the sorted positions spread out across $(0,1)$ without overlapping.}
\label{ch16:fig-order}
\end{figure}

\begin{paradoxbox}
\textbf{The German Tank Problem.} In World War II, the Allies wanted to estimate how many
tanks Germany was producing. Captured tanks had sequential serial numbers
$1,2,\dots,N$, and the Allies observed a random sample of $n$ of them; $N$ was unknown.
The largest serial number seen, $X_{(n)}$, is an order statistic, and naively it
\emph{underestimates} $N$ -- the true maximum is almost surely larger than any sample draw.
Order-statistic geometry fixes the bias: the $n$ observed numbers divide the range
$\{1,\dots,N\}$ into $n+1$ equal expected gaps, so the average gap between adjacent values
is $(N-n)/(n+1)$. The sample maximum therefore expects to reach $\E[X_{(n)}] = \frac{n}{n+1}(N+1)$.
Solving for $N$ yields the minimum-variance unbiased estimator:
\[
  \hat N = X_{(n)} + \frac{X_{(n)}}{n} - 1 = \frac{n+1}{n}X_{(n)} - 1.
\]
Statisticians using this beat the intelligence agencies' guesses dramatically -- postwar
German factory records revealed the mathematical formula was off by just a few percent
while espionage estimates were off by factors of two to three. A simple transformation
of an order statistic outperformed an international spy network.
\end{paradoxbox}

\section{The Probability Integral Transform}
\label{ch16:pit}

We now reach the most useful transformation of all. It comes from applying one very
special function $g$ -- a random variable's own CDF -- to that random variable.

\begin{intuitionbox}
The CDF $F$ measures \emph{how much probability lies to the left}. If you feed $X$ into its
own $F$, you get $U = F(X)$ -- the percentile rank of $X$. And percentile ranks, by
construction, are spread perfectly evenly: a quarter of the data has rank below $0.25$,
half has rank below $0.5$, and so on. That ``perfectly even'' is the definition of the
uniform distribution. So $F(X)$ is always $\Unif(0,1)$, no matter what $X$ was.
\end{intuitionbox}

\begin{theorem}[Probability Integral Transform]
\label{ch16:pit-thm}
Let $X$ be a continuous random variable with a strictly increasing CDF $F$. Then
\[
  U := F(X) \sim \Unif(0,1).
\]
Conversely, if $U \sim \Unif(0,1)$, then $X := F^{-1}(U)$ has CDF $F$.
\end{theorem}

\begin{proof}
For the forward direction, take $u \in (0,1)$. Because $F$ is strictly increasing it has an
inverse, and
\[
  \P(U \le u) = \P\big(F(X) \le u\big) = \P\big(X \le F^{-1}(u)\big)
              = F\big(F^{-1}(u)\big) = u.
\]
A random variable whose CDF is $\P(U\le u) = u$ on $(0,1)$ is exactly $\Unif(0,1)$.

For the converse, let $U \sim \Unif(0,1)$ and $X = F^{-1}(U)$. Then for any $x$,
\[
  \P(X \le x) = \P\big(F^{-1}(U) \le x\big) = \P\big(U \le F(x)\big) = F(x),
\]
using that $\P(U \le u) = u$. So $X$ has CDF $F$, as claimed.
\end{proof}

\begin{happeningbox}
\textbf{Equivalence with the uniform distribution.} The two halves of the theorem say
the same thing from opposite ends. The forward half flattens \emph{any} continuous
distribution into the uniform by viewing each value through its own percentile lens. The
backward half rebuilds \emph{any} distribution from the uniform by running the percentile
lens in reverse with $F^{-1}$, known as the \textbf{quantile function}. The quantile
function takes a probability rank $u \in (0,1)$ and returns the cutoff value $x$ below
which that fraction of probability sits. The uniform distribution on $(0,1)$ thus
provides a universal foundation: given uniform random numbers and an invertible CDF, one
can generate samples from any desired continuous distribution. This principle forms the
basis of stochastic simulation.
\end{happeningbox}

\begin{recoverybox}
\textbf{Pause and think.} Suppose $U \sim \Unif(0, 1)$ and you define $X = -\ln(1-U)$. What is
$\P(X > 2)$? Solve this in two ways: (1) using basic algebra on uniforms, and (2) using the
exponential distribution.
\emph{Check:} (1) $X > 2 \iff -\ln(1-U) > 2 \iff 1-U < e^{-2} \iff U > 1 - e^{-2}$, which has
probability $1 - (1 - e^{-2}) = e^{-2}$. (2) By the probability integral transform,
$X \sim \Expo(1)$, whose survival function is directly $\P(X > 2) = e^{-2}$. The two approaches
match instantaneously.
\end{recoverybox}

\begin{historybox}
The realization that a cumulative distribution function turns its own random variable into a
uniform crystallized during the 1930s. Statisticians like Karl Pearson and Ronald A.\ Fisher
needed a universal framework for ``goodness-of-fit'' tests -- checking whether observed
experimental data genuinely follows a proposed scientific model. Under Fisher's framework,
applying $F$ to a test statistic under the ``null hypothesis'' (the baseline assumption that no
special effect exists) produced a $p$-value that was guaranteed to be flatly uniform on $(0,1)$.
A controversy immediately broke out: Fisher viewed $p$-values as a continuous scale of evidentiary
surprise, while Jerzy Neyman and Egon Pearson insisted on rigid, pre-set error thresholds for
making binary decisions.

A decade later in 1946 at Los Alamos, Stanislaw Ulam and John von Neumann hit a computational wall.
They were trying to model neutron diffusion -- the chaotic scattering and chain-reaction fission
of neutrons inside fissile cores -- but the governing differential equations were far too
complex to solve on paper. Their only available tools were mechanical Marchant desk calculators,
cumbersome published tables of random digits, and the newly built ENIAC, an experimental
computer containing 18,000 glowing vacuum tubes.

Ulam realized they could bypass analytic integration by playing a game of chance on the computer:
simulate thousands of individual neutron trajectories step by step and tally the outcomes.
Nicholas Metropolis dubbed the technique the \emph{Monte Carlo method}, after the Monaco casino
where Ulam's uncle borrowed money to gamble. Inverting the CDF was their primary tool for turning
raw uniform numbers into physical velocities and collision angles. But this sparked a second
controversy: could a deterministic machine produce genuine randomness? Von Neumann famously
admitted that ``anyone who considers arithmetical methods of producing random digits is, of
course, in a state of sin.'' Critics worried that deterministic pseudorandom formulas would have
hidden repeating cycles that would corrupt the nuclear physics simulations.

Monte Carlo simulation soon became indispensable across scientific computing, Bayesian inference, and statistical
mechanics. Today, the foundational open challenge lies in high dimensions: when dealing with
millions of coupled variables -- such as pixels in image generation or tokens in language models --
the joint CDF is impossible to write down or invert directly. Developing efficient transformations
from simple Gaussian or uniform noise to intractable, high-dimensional target distributions
remains one of the most active research frontiers in machine learning.
\end{historybox}

\begin{applicationbox}
\textbf{Inverse-transform sampling.} To generate a sample from a distribution with CDF $F$,
when you can compute the quantile function $F^{-1}$:
\begin{enumerate}
  \item Draw $U \sim \Unif(0,1)$ (every language gives you this).
  \item Return $X = F^{-1}(U)$.
\end{enumerate}
By Theorem~\ref{ch16:pit-thm}, $X$ has exactly the CDF $F$. For the exponential, $F(x) =
1 - e^{-\lambda x}$ inverts to $F^{-1}(u) = -\tfrac{1}{\lambda}\ln(1-u)$, so
$X = -\tfrac{1}{\lambda}\ln(1-U)$ -- or, since $1-U$ is also uniform,
$X = -\tfrac{1}{\lambda}\ln U$ -- is an exponential sample. Two arithmetic operations turn
a uniform into an exponential. This same trick generates Cauchy ($\tan(\pi(U-\tfrac12))$),
logistic, Pareto, and any distribution whose quantile you can write down.
\end{applicationbox}

\begin{labbox}
Before reading on, derive an inverse-transform generator for the standard \textbf{logistic}
distribution, whose CDF is $F(x) = \dfrac{1}{1+e^{-x}}$ (the sigmoid you may know from
logistic regression). Solve $u = 1/(1+e^{-x})$ for $x$. You should find
$X = \ln\!\big(U/(1-U)\big)$ -- the \emph{log-odds} of a uniform. Notice that the sigmoid
and the log-odds (logit) are inverse functions, and that the probability integral
transform is what makes them a matched encoder/decoder pair.
\end{labbox}

\begin{figure}[H]
\centering
\begin{tikzpicture}
\begin{axis}[
  width=12cm, height=6cm,
  xlabel={$u$ (uniform input)}, ylabel={$x=F^{-1}(u)$ (output sample)},
  domain=0.01:0.99, samples=120,
  ymin=-0.2, ymax=4.6, xmin=0, xmax=1,
  legend pos=north west, legend cell align=left,
  axis lines=left, thick,
]
\addplot[densityred, very thick] {-ln(1-x)};
\addlegendentry{$F^{-1}(u)=-\ln(1-u)$, \ $\Expo(1)$}
\end{axis}
\end{tikzpicture}
\caption{Inverse-transform sampling for the $\Expo(1)$ distribution. Uniform inputs on the
horizontal axis are pushed through the quantile function $F^{-1}$. Note how the steep right
end stretches the rare large uniforms into the long exponential tail -- the curve's slope
is doing the ``paint relocation'' of Section~\ref{ch16:core}.}
\label{ch16:fig-pit}
\end{figure}

\begin{trapbox}
The probability integral transform as stated needs $F$ \emph{continuous} (so that $F(X)$
hits every value in $(0,1)$) and is cleanest when $F$ is strictly increasing (so $F^{-1}$
is an ordinary two-sided inverse function). For discrete or mixed distributions, $F$ has
jumps and flat stretches. In that case, one uses the \textbf{generalized inverse}:
\[
  F^{-1}(u) = \inf\set{x : F(x) \ge u},
\]
which identifies the smallest value $x$ where the cumulative probability staircase reaches or
exceeds percentile rank $u$. Inverse-transform sampling still works flawlessly with this
definition -- it is how a computer generates a Poisson, Geometric, or Bernoulli variable
from a single uniform draw. However, because $F(X)$ takes only discrete values, $F(X)$ is no
longer uniform when $X$ is discrete, so do not run the forward half blindly on discrete
distributions.
\end{trapbox}

\section{Box--Muller: Manufacturing Normals}
\label{ch16:box-muller}

Inverse-transform sampling hits one major snag: the normal distribution. Its CDF $\Phi$ has
\emph{no elementary closed form}, and neither does $\Phi^{-1}$. So the slick
``$F^{-1}(U)$'' recipe stalls exactly where we need it most -- the normal is the most
sampled distribution on Earth. In 1958, George Box and Mervin Muller introduced an elegant
solution by transforming \emph{two} uniforms simultaneously using polar coordinates in the plane.

\begin{intuitionbox}
The trick: a single normal is hard, but a \emph{pair} of independent normals is easy --
because their joint density depends only on the distance from the origin. Picture a
dartboard whose darts follow a 2D Gaussian: the cloud is perfectly circular. So instead of
sampling the two Cartesian coordinates, sample in \emph{polar} coordinates -- pick a random
angle (easy: uniform) and a random radius (also easy, via inverse transform on the radius's
nice distribution). Then convert back to $(x,y)$ and read off two independent normals.
\end{intuitionbox}

\begin{whybox}
\textbf{Why the radius is easy.} Let $X,Y \sim \Normal(0,1)$ independently. Their joint
density is
\[
  f(x,y) = \frac{1}{2\pi}\,e^{-(x^2+y^2)/2}.
\]
It depends on $(x,y)$ only through $r^2 = x^2+y^2$ -- the circular symmetry made visible.
Switch to polar coordinates $x = r\cos\theta$, $y=r\sin\theta$. The Jacobian of the polar
map is $r$ (the standard area element $\dd x\,\dd y = r\,\dd r\,\dd\theta$), so the joint
density of $(R,\Theta)$ is
\[
  f_{R,\Theta}(r,\theta) = \frac{1}{2\pi}\,e^{-r^2/2}\cdot r
  = \underbrace{\frac{1}{2\pi}}_{\Theta}\;\cdot\;\underbrace{r\,e^{-r^2/2}}_{R},
  \qquad 0<r,\ 0\le\theta<2\pi.
\]
This \emph{factors}: $\Theta \sim \Unif(0,2\pi)$ and, independently, $R$ has density
$r\,e^{-r^2/2}$. The angle is uniform -- of course, by symmetry -- and the radius has a
density we can integrate in closed form, unlike $\Phi$.
\end{whybox}

\begin{figure}[H]
\centering
\begin{tikzpicture}[scale=1.0]
\draw[->,gray] (-2.6,0) -- (2.8,0) node[right, font=\small]{$x=Z_1$};
\draw[->,gray] (0,-2.6) -- (0,2.8) node[above, font=\small]{$y=Z_2$};

\foreach \a/\rr in {75/1.8, 80/1.0, 110/2.0, 130/1.5, 160/0.8, 190/1.6, 210/1.1, 240/1.8, 260/0.9, 285/1.5, 310/1.0, 335/1.7, 350/0.8}{
  \fill[sourceblue!55] (\a:\rr) circle (1.6pt);
}

\def\ang{38}\def\rad{2.4}
\draw[targetgreen, very thick] (0,0) -- (\ang:\rad)
  node[pos=0.58, above left=-1pt, font=\footnotesize] {$R=\sqrt{-2\ln U_1}$};
\fill[densityred] (\ang:\rad) circle (2.4pt);
\draw[densityred, thick] (0.75,0) arc (0:\ang:0.75);
\node[densityred, font=\footnotesize, anchor=west] at (16:0.85) {$\Theta=2\pi U_2$};
\node[densityred, right=2pt, font=\footnotesize] at (\ang:\rad) {$(Z_1,Z_2)$};
\end{tikzpicture}
\caption{Box--Muller geometry. A pair of independent standard normals forms a rotationally
symmetric Gaussian cloud in the plane. Rather than sampling the Cartesian coordinates directly, draw
a uniform angle $\Theta=2\pi U_2$ and an independent radius $R=\sqrt{-2\ln U_1}$, then project back
via $(Z_1,Z_2)=(R\cos\Theta,\,R\sin\Theta)$.}
\label{ch16:fig-boxmuller}
\end{figure}

Now sample $R$ by inverse transform. Its CDF is
\[
  F_R(r) = \int_0^r s\,e^{-s^2/2}\,\dd s = 1 - e^{-r^2/2},
\]
so $S := R^2$ has $\P(R^2 \le s) = 1 - e^{-s/2}$ -- that is, $R^2 \sim \Expo(\tfrac12)$,
which we already know how to sample. Setting $1 - e^{-r^2/2} = U_1$ and solving,
$R = \sqrt{-2\ln(1-U_1)}$, or equivalently $R = \sqrt{-2\ln U_1}$. Convert back:

\begin{formalbox}
\textbf{Box--Muller transform} (Figure~\ref{ch16:fig-boxmuller})\textbf{.} Let $U_1, U_2$ be independent $\Unif(0,1)$. Define
\[
  R = \sqrt{-2\ln U_1}, \qquad \Theta = 2\pi U_2,
\]
and set
\[
  \boxed{\;Z_1 = R\cos\Theta = \sqrt{-2\ln U_1}\,\cos(2\pi U_2), \qquad
          Z_2 = R\sin\Theta = \sqrt{-2\ln U_1}\,\sin(2\pi U_2).\;}
\]
Then $Z_1$ and $Z_2$ are \emph{independent} $\Normal(0,1)$ random variables. Two uniforms
in, two independent normals out.
\end{formalbox}

\begin{happeningbox}
What just happened is the multivariate change-of-variables theorem
(Section~\ref{ch16:jacobian}) run \emph{in reverse} and read as an algorithm. We could not
invert $\Phi$ in one dimension, so we lifted to two dimensions where the Gaussian's
rotational symmetry hands us a factorization the one-dimensional problem hid. The Jacobian
factor $r$ from polar coordinates makes this possible: it converts $e^{-r^2/2}$ into
$r\,e^{-r^2/2}$, which has an elementary antiderivative $-e^{-r^2/2}$.
\end{happeningbox}

\begin{example}[Box--Muller by hand]
\label{ch16:bm-example}
Suppose your random-number generator hands you $U_1 = 0.135$ and $U_2 = 0.250$. Then
\[
  R = \sqrt{-2\ln 0.135} = \sqrt{-2(-2.002)} = \sqrt{4.005} \approx 2.001,
  \qquad \Theta = 2\pi(0.250) = \tfrac{\pi}{2}.
\]
So $Z_1 = 2.001\cos(\pi/2) \approx 0$ and $Z_2 = 2.001\sin(\pi/2) \approx 2.001$. We have
produced two standard-normal draws, $0.00$ and $2.00$, from two uniforms using a
log, a square root, and a sine/cosine. (Here $U_2 = 0.25$ landed the angle exactly on the
$y$-axis, which is why $Z_1$ came out essentially zero -- a coincidence of the chosen
numbers, not a bias.)
\end{example}

\begin{applicationbox}
\textbf{Where Box--Muller lives today.} Every time a deep-learning framework initializes a
weight matrix with Gaussian noise, every time a Monte Carlo option pricer simulates a
stock path, every time a physics engine adds thermal jitter, normal samples are being
generated -- and for decades Box--Muller (or its rejection-based cousin, the
\emph{polar/Marsaglia} method, which avoids the trig calls) did the job. Modern libraries
often use the even faster Ziggurat algorithm, but Box--Muller remains the one everyone
\emph{derives}, because it shows so cleanly how a clever transformation turns the cheap and
universal uniform into the expensive and indispensable normal.
\end{applicationbox}

\section{A Unifying View}
\label{ch16:unify}

\begin{aestheticbox}
All three methods express the same core transformation principle in different mathematical forms.

The \textbf{CDF method} tracks accumulated probability directly and works for any
distribution without stretch factors. The \textbf{Jacobian method} is the CDF method differentiated --
it trades global accounting for a local stretch factor, and it is fastest when $g$ is
smooth and invertible. \textbf{Order statistics} are a multivariate transformation
(sorting) handled by either method, plus a combinatorial count of arrangements. And the
\textbf{probability integral transform} is the special case where the transforming
function is the CDF itself, collapsing every distribution onto the uniform and back.

The single thread running through all of it is \emph{conservation of probability}:
probability is a fixed total of one unit, and a transformation can only \emph{move} it,
never make or destroy it. Master that one principle and you never need to memorize a
transformation formula again -- you can re-derive any of them from ``where did the
probability go?''
\end{aestheticbox}

\subsection*{Looking ahead: sums and convolutions}
The most important transformation in probability is the sum of two independent
random variables, $S = X + Y$. Applying the bivariate change of variables
$(S, T) = (X+Y, X)$ and integrating out $T$ yields the continuous \textbf{convolution formula},
\[
  f_S(s) = \int_{-\infty}^{\infty} f_X(t)\,f_Y(s-t)\,\dd t.
\]
Chapter~17 develops convolutions for discrete and continuous variables, exploring
how distributions behave under addition and leading directly toward the Central Limit Theorem.

\section*{Exercises}
\addcontentsline{toc}{section}{Exercises}

\begin{exercise}[CDF method, monotone]
\label{ch16:ex-cube}
Let $X \sim \Unif(0,1)$ and $Y = X^3$. Find $F_Y$ and $f_Y$ by the CDF method, then
re-derive $f_Y$ using the Jacobian formula and check they agree. Where does the density
blow up, and why?
\end{exercise}

\begin{exercise}[Folding a symmetric variable]
\label{ch16:ex-abs}
Let $X \sim \Normal(0,1)$ and $Y = \abs{X}$ (the ``half-normal''). Show that
$f_Y(y) = \sqrt{2/\pi}\,e^{-y^2/2}$ for $y>0$, and explain the factor of $2$ in terms of
preimages.
\end{exercise}

\begin{exercise}[Exponential to uniform and back]
\label{ch16:ex-expo-pit}
Let $X \sim \Expo(\lambda)$. Verify directly that $U = 1 - e^{-\lambda X}$ is $\Unif(0,1)$,
illustrating the forward probability integral transform. Then start from
$U \sim \Unif(0,1)$ and recover an exponential, confirming the converse.
\end{exercise}

\begin{exercise}[Inverse-transform design]
\label{ch16:ex-pareto}
The Pareto distribution has CDF $F(x) = 1 - (x_m/x)^{\alpha}$ for $x \ge x_m > 0$ and
$\alpha>0$. Derive the inverse-transform sampler $X = F^{-1}(U)$, and write a one-line
formula that turns a uniform draw into a Pareto draw.
\end{exercise}

\begin{exercise}[Order statistics computation]
\label{ch16:ex-order}
Five sensors have independent $\Unif(0,1)$ failure thresholds. (a) Find the density of the
\emph{second-smallest} threshold $X_{(2)}$ and identify it as a Beta distribution. (b)
Compute $\E[X_{(2)}]$. (c) Find the probability that the largest threshold exceeds $0.9$.
\end{exercise}

\begin{exercise}[Minimum of exponentials]
\label{ch16:ex-min-expo}
Let $X_1,\dots,X_n$ be independent with $X_i \sim \Expo(\lambda_i)$ (possibly different
rates). Show that $\min_i X_i \sim \Expo\!\big(\sum_i \lambda_i\big)$, and that the
probability that $X_j$ is the minimum equals $\lambda_j / \sum_i \lambda_i$. (Hint: use the
survival function for the first part; condition on $X_j$ and integrate for the second.)
\end{exercise}

\begin{exercise}[A two-variable Jacobian]
\label{ch16:ex-jacobian2}
Let $X, Y$ be independent $\Expo(1)$. Define $S = X + Y$ and $T = X/(X+Y)$. Using the
two-dimensional change-of-variables formula, show that $S \sim \Gam(2,1)$ and
$T \sim \Unif(0,1)$, and that $S$ and $T$ are independent. (This is a famous and surprising
decoupling of ``total'' from ``proportion.'')
\end{exercise}

\begin{exercise}[Box--Muller verification]
\label{ch16:ex-bm}
Carry out the Box--Muller change of variables in full: starting from the joint density of
$(U_1, U_2)$, with $(Z_1, Z_2)$ as the map of $(U_1, U_2)$, compute the $2\times2$ Jacobian determinant and
confirm that the resulting joint density of $(Z_1,Z_2)$ is
$\tfrac{1}{2\pi}e^{-(z_1^2+z_2^2)/2}$, hence two independent standard normals.
\end{exercise}

\begin{exercise}[Conceptual: why the absolute value?]
\label{ch16:ex-absval}
Explain, without computation, why the change-of-variables formula has $\abs{g'(x)}$ rather
than $g'(x)$. What would go wrong for a strictly decreasing $g$ if the absolute value were
omitted? Relate your answer to ``densities are nonnegative'' and ``probabilities are
nonnegative.''
\end{exercise}

\begin{exercise}[Proof: PIT for general continuous $F$]
\label{ch16:ex-pit-proof}
Prove the forward probability integral transform without assuming $F$ is strictly
increasing -- only continuous. (Hint: define the generalized inverse
$F^{-1}(u) = \inf\set{x : F(x)\ge u}$ and show $\P(F(X)\le u) = u$ by relating the event
$\set{F(X)\le u}$ to $\set{X \le F^{-1}(u)}$.)
\end{exercise}

\begin{exercise}[Median of an odd sample]
\label{ch16:ex-median}
For an i.i.d.\ sample of size $n = 2m+1$ from $\Unif(0,1)$, the sample median is
$X_{(m+1)}$. Write its density, identify the Beta distribution, and argue from symmetry
that its mean is exactly $\tfrac12$. How does its variance behave as $n \to \infty$, and
what does that say about the sample median as an estimator?
\end{exercise}

\begin{exercise}[Simulating a triangular distribution]
\label{ch16:ex-triangular}
The symmetric triangular density on $(0,1)$ is $f(x) = 4x$ for $0<x\le\tfrac12$ and
$f(x)=4(1-x)$ for $\tfrac12<x<1$. (a) Find its CDF. (b) Derive the piecewise
inverse-transform sampler. (c) Show that the same distribution arises as the average
$(U_1+U_2)/2$ of two independent uniforms, giving a second, slicker simulation recipe.
\end{exercise}
